\newpage
\section{Missió 1: Compartir és estimar}\label{sec:mission1}

\begin{wrapfigure}{L}{0.23\linewidth}
\vspace{-0.5cm}\includegraphics[width=\linewidth]{abacus}
\end{wrapfigure}
%
Fa uns 75 anys, es va enviar una expedició d'alt secret a través de l'espai amb l'objectiu de colonitzar un altre planeta. L'expedició va arribar sana i estàlvia a la seva destinació, però l'aterratge va trencar els instruments de comunicació i els va deixar incomunicats. Ets un colon de tercera generació, futur enginyer, i aquesta és la primera missió que t'assigna l'assemblea:
compartir la nova i lluent \mbox{calculadora}. \textbf{Si us plau, llegeix-ho tot amb \mbox{atenció}}.\\

\noindent
\begin{remark}
\textbf{Important}: Es recomana que cada estudiant publiqui un \textbf{informe formal} (en PDF) al ``fòrum d'intercanvi de la Missió 1'' del Campus Virtual abans de la següent sessió de 2h. Després, comptem amb tu per publicar una \textbf{carta formal de resposta amb comentaris} (també en PDF) sobre almenys una altra solució, abans de la següent sessió d'1h.
\end{remark}

\subsection*{Detalls de l'informe}

Ara mateix, la colònia consta de 3 bases separades (3 edificis). Aquestes bases tenen 3, 200 i 5000 dispositius electrònics, respectivament. La base amb 3 dispositius s'ha construït específicament per allotjar una calculadora nova, potent i que consumeix molta energia (és un dels 3 dispositius). Tothom a la colònia vol tenir accés a la calculadora, però cap dels dispositius té capacitats de xarxa \textit{encara}. El teu objectiu és proposar un disseny de xarxa convincent que connecti tothom i els deixi a tots satisfets.\\

\noindent
Prepara un informe breu amb el teu disseny, amb propostes de solució clares i completes per als punts següents:

\begin{enumerate}[itemsep=0.1cm]
\item Com es travessa la \textbf{barrera analògic-digital} per permetre la comunicació digital? Sigues concret en les tecnologies que proposis.
\item Com es connecten físicament els dispositius? Proporciona un diagrama amb la \textbf{topologia} exacta.
\item Quin \textbf{sistema d'adreçament} s'utilitzarà per identificar de manera única cada dispositiu? Dona dos exemples d'adreça, un dels quals ha de ser el de la nova calculadora.
\item Quin \textbf{format de paquet} faran servir els dispositius? Proporciona un diagrama concret del format i els bytes d'un paquet d'exemple complet.
\end{enumerate}


\subsection*{Detalls dels comentaris}

Després de publicar la teva solució al fòrum, busca'n una altra i \textbf{reclama-la} responent-hi amb un missatge breu de ``Hi faré comentaris'' (una carta de comentaris per solució).\\

\noindent
Després prepara una carta formal en PDF i publica-la al mateix fil quan estigui llesta (trobaràs una carta de comentaris d'exemple al Campus Virtual). En aquesta carta, inclou els teus comentaris (tant positius com negatius), explicant qualsevol cosa que consideris rellevant o que creguis que pot ajudar aquella solució.  Assegura't de tractar els punts següents:

\begin{enumerate}[itemsep=0.1cm]
\item Hi veus alguna falla important o algun punt incomplet que impediria que la proposta es fes realitat?
\item Si tornessis a fer el teu informe, adoptaries alguna de les propostes d'aquesta solució?
\item Has pogut interpretar el contingut del missatge inclòs a l'informe? Inclou els bytes d'una possible resposta, basant-te en el que has entès d'aquell informe.
\end{enumerate}

